\documentclass[11pt,a4paper]{article}
\usepackage[T1]{fontenc}
\usepackage[utf8]{inputenc}
\usepackage{lmodern,microtype}
\usepackage[margin=25mm]{geometry}
\usepackage{amsmath,amssymb,amsthm,mathtools,mathrsfs,booktabs}
\usepackage[hyphens]{url}
\usepackage[hidelinks]{hyperref}
\setlength{\emergencystretch}{2em}
\newcommand{\E}{\mathbb E}
\newcommand{\Pp}{\mathbb P}
\newcommand{\R}{\mathbb R}
\newcommand{\CE}{C_{\mathrm E}}
\newtheorem{theorem}{Theorem}[section]
\newtheorem{lemma}[theorem]{Lemma}
\title{Initial acquisition and sharper count comparison\\
for fixed-target Berry--Esseen excess}
\author{Bartosz Ko{\l}odziejek \and Kacper Rodziewicz}
\date{27 October 2026}
\begin{document}
\maketitle
\begin{abstract}
We continue the fixed-target approach to the sharp iid Berry--Esseen
problem. The first support and moment bounds use fixed-target contact,
an external structural inequality and the sharp predecessor bound.
No upper bound on the actual score is needed. A moment-sensitive
Gaussian-weight condition then permits a small
complete scalar cover of the initial moment interval. Retaining the phase
cost improves its defect bound, although the remaining bound is too large
for the next general contraction step. Separately, disjoint crossing events
and a fourth-moment estimate reduce the cost in the conditional-count
comparison. A combined moment response gives a further conditional
improvement on one central domain. We give the new proofs and distinguish
them from the retained Fourier estimates, domain-acquisition premises and
finite-base obligations. No complete sharp upper bound is claimed.
\end{abstract}

\paragraph{How the pieces fit.}
The earlier companion establishes attainment, contact and conservation for
the excess $\sqrt n(F_n-\Phi)-\CE\rho$. Here we first make its initial
entry quantitative, then calculate what that entry actually yields. The
count-crossing argument and the common moment response are separate local
improvements; neither can be used before its stated domain is acquired.
The proof of each implication is given below. The low- and high-frequency
inequalities named in Section~\ref{v5entry:raw} remain explicit analytic
inputs, as in the full manuscript. Exact evaluation of their formulas does
not independently prove Fourier inversion or their uniform remainder bounds.

The external affine inequality used for initial entry is Shevtsova's
2013 structural bound, valid for every sample size; see
\emph{Informatics and its Applications} \textbf{7}(1), 124--125,
\url{https://www.mathnet.ru/eng/ia252}. It is applied after variance
normalization, where the third-moment sum is $\rho/\sqrt n$ and the sum of
cubed standard deviations is $1/\sqrt n$. Its uniform bound also applies
to closed CDF values by taking one-sided limits.

% Inputtable proof note. Wrapper: entry-main.tex.
\section{Entry from fixed-target excess}
\label{v5entry:section}

This note isolates the first acquisition step. It does not assert a complete
local exclusion or a sharp Berry--Esseen theorem. Let $P$ be a centered,
unit-variance law, $S_j$ a sum of $j$ independent copies, and set
\[
 \rho=\E|X|^3,\quad \beta=\E[X|X|],\quad \gamma=\E X^3,
 \quad c=\CE=\frac{3+\sqrt{10}}{6\sqrt{2\pi}}.
\]
The threshold $H$ always has the units of the raw sum. A fixed-target
maximizer maximizes
\[
 a(P,H)=\sqrt n\{\Pp(S_n\le H)-\Phi(H/\sqrt n)\}-c\rho.
\]
We use the attainment and contact theorem in the companion fixed-target
note. In particular its closed-support contact
equality, not ratio optimality, is the variational premise here. We consider
a maximizer with $a\ge0$. The actual normalized score is $c+a/\rho$;
no upper bound on it is assumed.

\begin{theorem}[Initial moment and support entry]\label{v5entry:initial}
Assume $n\ge42$, $C_{n-1}\le c$, and the universal structural inequality
\begin{equation}\label{v5entry:structural}
 \sqrt n\,\Delta_n(P)\le\frac{3031}{10000}
                  \left(\rho+\frac{323}{500}\right)
\end{equation}
for standardized laws with finite third moment. Then a fixed-target
maximizer with $a\ge0$ satisfies
\begin{equation}\label{v5entry:entry}
 1\le\rho<\frac{7361}{4000}<\frac{1841}{1000},\qquad
 |\beta|<\frac{321}{500},\qquad
 \operatorname{supp}P\subset(-96/25,96/25).
\end{equation}
The theorem is uniform in all real $H$ and in both signs of a support
point. It requires neither $C_{n-2}$ nor an excess bound of order $1/n$.
\end{theorem}

Inequality \eqref{v5entry:structural} is an external structural input,
recorded with its Shevtsova attribution in the manuscript introduction.
It is not proved by the exact arithmetic below. Only this one of the
two affine structural bounds is needed for the displayed moment radius.

\begin{proof}
Since $c\rho\le\sqrt n\Delta_n(P)$ and $c>819/2000$, elementary division gives
\[
 \rho<\frac{(3031/10000)(323/500)}{819/2000-3031/10000}
       =\frac{7361}{4000}.
\]
The other affine bound $0.3322(\rho+0.429)$ alone would give the larger
number $712569/386500$; it must not be used to claim $7361/4000$.
Put $m=\E X_+=\E X_->0$ and $a_\pm=\E X_\pm^2=(1\pm\beta)/2>0$.
Cauchy on each sign class gives $m\rho\ge a_+^2+a_-^2$ and
$m^2\le a_+a_-$. Hence
\[
 \rho\ge\frac{1+\beta^2}{\sqrt{1-\beta^2}}.
\]
This function increases with $|\beta|$. For $r=1841/1000$, $b=321/500$,
\[
 (1+b^2)^2-r^2(1-b^2)=\frac{93288849}{50000000000}>0,
\]
which proves the second assertion without assuming bounded support.

Write $v=n-1$, $s=\sqrt{n/v}$, $F_v(x)=\Phi(x/\sqrt v)$ and
$e_n=n(s-1)-1/2$. Closed-support contact, the upper predecessor
bound $F_{n-1,P}(H-t)\le F_v(H-t)+c\rho/\sqrt v$, and
$\sqrt n(F_{n,P}(H)-\Phi(H/\sqrt n))=c\rho+a$ give
\begin{equation}\label{v5entry:contact}
 0\le H_n(H,t)+c\{-|t|^3+\tfrac32\rho t^2+3\beta t+\rho e_n\}-na
 \quad(t\in\operatorname{supp}P),
\end{equation}
where exactly
\[
 \frac{H_n(H,t)}{n^{3/2}}
 =F_v(H-t)-F_{v+1}(H)+tF'_{v+1}(H)
                   -\frac{t^2-1}{2}F''_{v+1}(H).
\]
The extra $-na$ can be discarded in this upper bound. This is the precise
reason the predecessor assumption is $C_{n-1}\le c$, not a comparison
with the actual score. No continuity assumption on the distribution
function at $H-t$ is inserted.

Rationalization gives
\[
 ne_n=\frac{s^2(2s+1)}{2(s+1)^2},\qquad
 s<81/80=:S,\qquad
 ne_n<\frac{793881}{2073680}<\frac{48}{125}.
\]
The first expression increases for $s\ge1$, as its logarithmic derivative
is $2/s+2/(2s+1)-2/(s+1)>0$. Thus $e_n<8/875$.

For completeness we rederive the Gaussian estimate used at the radius.
The heat identity $\partial_vF_v=F_v''/2$ permits spatial expansion of
$F_v(H-t)$ to order four, variance expansion of $F_{v+1}$ to order two,
and variance expansion of $F'_{v+1},F''_{v+1}$ to order one. The cubic
coefficient is $(-t^3+3t)/6$, and the quartic coefficient is
$(t^4-6t^2+3)/24$. Therefore, for $t=\pm u$, $z_v=H/\sqrt v$,
\begin{equation}\label{v5entry:heat}
 H_n(H,\pm u)/s^3
 =\pm A(1-z_v^2)\phi(z_v)+B\phi'''(z_v)+\mathcal E,
\end{equation}
where $A=(u^3-3u)/6$, $H_4=u^4-6u^2+3$ and $B=H_4/(24\sqrt v)$.
The four error bounds after division by $s^3$ are respectively
\[
 \frac{u^5\|\phi^{(4)}\|_\infty}{120v},\quad
 \frac{\|\phi^{(5)}\|_\infty}{48v^{3/2}},\quad
 \frac{u\|\phi^{(4)}\|_\infty}{8v},\quad
 \frac{(u^2-1)\|\phi^{(5)}\|_\infty}{16v^{3/2}}.
\]
Gaussian Fourier inversion gives norms at most $3\phi(0)<6/5$ and
$8/\pi<51/20$. Since $v\ge41$ and $\sqrt v>6$, for $u\ge5/2$ this gives
\[
 |\mathcal E|<E(u):=u^5/4100+3u/820+17(3u^2-2)/78720.
\]
These errors are uniform in $H$; they use no higher moment of $P$.

If $A,B>0$ and $B\le A/4$, the leading part of \eqref{v5entry:heat}
is at most $\phi(0)(A+9B^2/(4A))$ for both signs. Indeed, for the
positive sign and $|z_v|\le1$, it is at most
$\phi(0)(A-Az_v^2+3B|z_v|)$, and complete the square. Outside that
interval the first term is nonpositive. For the negative sign the first
term is at most $2e^{-3/2}A\phi(0)<A\phi(0)/2$. In both cases use
$\|\phi'''\|_\infty<2\phi(0)$: on $[0,\sqrt3]$ the polynomial
$3x-x^3$ is at most two, and beyond $\sqrt3$ use
$x^3e^{-x^2/2}\le3\sqrt3e^{-3/2}<2$.

For frozen $r,b$ above define
\begin{align*}
 Q(u)&=-u^3+\tfrac32ru^2+3bu+(8/875)r,\\
 P(u)&=\tfrac25 S^3\left(A^2+\frac{H_4^2}{256\cdot41}\right)
                         +S^3E(u)A+\frac{819}{2000}Q(u)A.
\end{align*}
At $R=96/25<6$ we have $A>0$, $0<H_4<36A$, hence $B<A/4$.
The sign $Q(R)<0$ is checked before replacing $c$ by its lower endpoint.
The upper bound in \eqref{v5entry:contact}, multiplied by $A$, is then
strictly below $P(R)$, and exact arithmetic gives
\[
 -P(R)=\frac{2214258283536204102723}{10250000000000000000000}>\frac15.
\]
To cover each entire exterior ray, differentiation of the exact Gaussian
expression gives
\[
 |\partial_tH_n(H,t)|<\tfrac25(t^2/2+1/2+e_n).
\]
For verification put $y=(H-t)/\sqrt n$, $z=H/\sqrt n$; the derivative is
$n\phi(z)+Ht\phi(z)-ns\phi(sy)$. The second-order Taylor error between
$\phi(y)$ and $\phi(z)$ costs $\phi(0)t^2/2$, while
$\partial_s[s\phi(sy)]=-\phi''(sy)$ costs $\phi(0)n(s-1)$.
Define $L(u)=\tfrac25(u^2/2+1/2+8/875)+(819/2000)Q'(u)$.
At $R$, the six numbers $Q,Q',Q'',L,L',L''$ are strictly negative.
Since $Q'''=-6$ and $L''=2/5-6(819/2000)<0$, all their required signs
persist for $u\ge R$. Thus both radial upper bounds decrease from a
strictly negative boundary value. The contact inequality excludes every
support point with $|t|\ge R$.
\end{proof}

\section{Moments and an actual high-frequency event}
\label{v5entry:moments}
Higher moments are used only after Theorem~\ref{v5entry:initial}.
For $x=|X|$ the nonnegative measure
$d\mu_0=(x-1)^2(x+1/2)\,dP$ has mass $\rho-1$. Direct expansion gives
\begin{align*}
 M_4-1&=\int\frac{(x+1)^2}{x+1/2}\,d\mu_0,\\
 M_5-1&=\int\frac{x^3+2x^2+3x+3/2}{x+1/2}\,d\mu_0,
\end{align*}
where $M_j=\E|X|^j$. Both ratios increase for $x\ge0$: they equal
$x+3/2+1/[4(x+1/2)]$ and
$x^2+3x/2+9/4+3/[8(x+1/2)]$, whose derivatives are nonnegative.
Writing their values at $R=96/25$ as $k_4(R),k_5(R)$, we obtain
\begin{equation}\label{v5entry:mom}
 M_4\le\min(R\rho,1+k_4(R)(\rho-1)),\quad
 M_5\le\min(RM_4,1+k_5(R)(\rho-1)),\quad M_6\le RM_5.
\end{equation}
Also $\gamma^2\le M_4-1$, by expanding
$\E(X^2-1-\gamma X)^2\ge0$.

There is no need to assume a central support gap before the first Fourier
bound. For independent copies let $E=\{(X-X')^2\ge9/25\}$ and let
$M$ be any valid fourth-moment ceiling. Then
\begin{equation}\label{v5entry:event}
 \Pp(E)\ge\alpha:=\frac{(2-9/25)^2}{2M+6}>0,\quad
 |X-X'|\ge3/5\text{ on }E,\quad
 \E[(X-X')^2\mid E]\le2/\alpha.
\end{equation}
Indeed $\E[(X-X')^2;E]\ge2-9/25$ and
$\E(X-X')^4=2M_4+6\le2M+6$, so Cauchy proves the mass bound.
The conditional second-moment bound follows from its unconditional value
two. This is an event under one actual iid marginal, not a replacement
characteristic function.

\section{A sharper first license}
\label{v5entry:license-section}
The retained fixed-target conserved identity is
\begin{equation}\label{v5entry:master}
 (n+H^2)a+nc\rho-\frac{n\phi(H/\sqrt n)}6(\gamma+3J)
                  +n^{3/2}\mathcal L=\mathcal R_c(H).
\end{equation}
Put $\tau(t)=\E[X\mathbf1_{X\ge t}]$ and $\nu=\mathcal L(H-S_{n-1})$.
Over complementary support gaps, define
\[
 J=\sum\tau_{u,v}(v-u)^2,\qquad
 \mathcal L=\sum\tau_{u,v}\int_{[u,v)}(t-u)\,d\nu(t)\ge0.
\]
The half-open interval includes the correct endpoint atoms. The elementary
gap inequality gives
$D=(3+\sqrt{10})\rho-\gamma-3J\ge0$.
In particular the right side of \eqref{v5entry:master} is nonnegative
when $a\ge0$, since it can also be written as
\[
 (n+H^2)a+nc\rho(1-e^{-w})+n\phi(0)e^{-w}D/6+n^{3/2}\mathcal L,
 \qquad w=H^2/(2n).
\]

\begin{lemma}[Moment-cell license]\label{v5entry:license}
Under \eqref{v5entry:master}, with $D,\mathcal L\ge0$, suppose
$n\ge N$, $\rho\in[l,u]$, $a\ge0$, and
$0<c_-\le c$. If a bound valid for every real $H$ gives
\begin{equation}\label{v5entry:gate}
 \mathcal R_c(H)\le T<Nc_-l,
\end{equation}
then either $T<0$ and this cell is excluded, or
\begin{gather}
 \gamma+3J>0,\qquad
 e^{-H^2/(2n)}\ge1-\frac{T}{Nc_-l}>0,\label{v5entry:density}\\
 (n+H^2)a+\frac{n\phi(0)}6D+n^{3/2}\mathcal L
                      \le\mathcal R_c(H)\le T.\label{v5entry:unweighted}
\end{gather}
Thus the first gate need not be the stronger uniform condition $T<2N/5$.
\end{lemma}
\begin{proof}
For a nonexcluded cell $T\ge0$. Divide \eqref{v5entry:master} by $n$:
\[
 \frac{\phi(H/\sqrt n)}6(\gamma+3J)
 =(1+H^2/n)a+c\rho+\sqrt n\mathcal L-\mathcal R_c/n
 \ge c\rho-T/n>0.
\]
This proves the sign. Since $D\ge0$,
$\gamma+3J\le6c\rho/\phi(0)$; hence
$c\rho e^{-w}\ge c\rho-T/n$, proving \eqref{v5entry:density}.
Finally rewrite the conserved expression as the left side of
\eqref{v5entry:unweighted}, plus
$n(1-e^{-w})(c\rho-\phi(0)D/6)$.
Its last coefficient is $\phi(0)(\gamma+3J)/6>0$, so it can be dropped.
\end{proof}

The positivity argument uses no small-defect premise. The input $D\ge0$
is used for the density and unweighted-resource conclusions, not for the
first sign deduction. Thresholds of either sign and $H=0$ are included.
If only $T\le Nc_-l$ is proved, the strict sign and positive density
floor do not follow. The strict gate is a genuine acceptance condition.

\begin{lemma}[First phase price without a small defect]\label{v5entry:phase}
Under Lemma~\ref{v5entry:license}, put
\[
 \mathcal U=(n+H^2)a+n\phi(0)D/6+n^{3/2}\mathcal L,
 \qquad q_{\rm ph}=\tfrac12(c_-l-T/N)>0.
\]
Then $\mathcal U\le\mathcal R_c(H)-q_{\rm ph}H^2$.
In particular a complete polynomial bound
$\mathcal R_c(H)\le A+F|H|-qH^2$ with $F\ge0,q>0$ gives
\[
 \mathcal U\le A+\frac{F^2}{4(q+q_{\rm ph})}.
\]
This improved bound is on $\mathcal U$, not on the raw $\mathcal R_c$.
\end{lemma}
\begin{proof}
Let $M_c=c\rho-\phi(0)D/6$ and $w=H^2/(2n)$.
The exact identity and $\mathcal R_c\le T$ imply
$e^{-w}M_c\ge c\rho-T/n>0$. Therefore
\[
 n(1-e^{-w})M_c=n(e^w-1)e^{-w}M_c
 \ge\frac{H^2}{2}(c\rho-T/n)\ge q_{\rm ph}H^2.
\]
Here $e^w-1\ge w$ for $w\ge0$. Subtract this term from the exact
decomposition $\mathcal R_c=\mathcal U+n(1-e^{-w})M_c$, then complete
the square. There is no premise $D<2$, and no multiplication of separate
lower estimates on $M_c$ and $e^{-w}$.
\end{proof}

\begin{lemma}[Normalized first-cover consequence]\label{v5entry:ten}
Under the preceding hypotheses, suppose a separate complete enclosure
proves $\mathcal R_c(H)/\rho<10$ for all real $H$. For $n\ge42$,
\[
 e^{-H^2/(2n)}>\frac{7199}{17199}>\frac{209}{500},\qquad
 \frac D\rho<\frac{25000}{6979},\qquad
 \frac a\rho<\frac{10}{n}.
\]
In particular the first license is discharged. The actual score satisfies
$c+a/\rho<c+10/n$, but this is an output, not an entry premise.
\end{lemma}
\begin{proof}
Use the pointwise version of the preceding proof with $T=10\rho$:
$10<42(819/2000)<nc$. The density floor is strictly larger than
$1-10/(42c)>1-10/[42(819/2000)]=7199/17199$.
From \eqref{v5entry:unweighted} and $\phi(0)>997/2500$,
$D/\rho<60/(42\phi(0))<25000/6979$.
The nonnegative term $(n+H^2)a$ gives the final bound.
\end{proof}

\section{The first scalar enclosure and its remaining gate}
\label{v5entry:raw}
This section states exactly what a bounded arithmetic cover must supply.
It proves the algebraic reduction, rather than treating an archived PASS
as an analytic premise. Fix a closed cell $[l,u]\subset[1,7361/4000]$.
Choose outward ceilings $M,g,E_5,b$ for $M_4,|\gamma|,M_5,|\beta|$ using
\eqref{v5entry:mom}, $|\gamma|\le\rho$, Pearson, and the sign-moment
inequality. In particular $M<11.1$, $g<1.85$ and $M_6\le RE_5$.

Write $J_1=\sum\tau_{a,b}\int_a^b(t-a)t\,dt$ and
$J_2=\sum\tau_{a,b}\int_a^b(t-a)t|t|\,dt$. Set
\begin{align*}
 K_1&=\E[X|X|^3]-3\beta-3\rho\gamma/2,\\
 K_q&=M_5/4-\rho M_4/2-3\beta\gamma/2+\rho/2,
 \qquad U=J_1+(M_4-3)/6.
\end{align*}
The exact raw remainder in \eqref{v5entry:master} is
\[
 \mathcal R_c=c(3\beta J/2+3\rho J_1-3J_2-K_q)
       +H(cK_1+\phi U)-H^2(c\rho+\gamma\phi/2)+E_n.
\]
The fixed target $c$ is the coefficient in every term. Its enclosing
interval $[819/2000,20487/50000]$ is not a bound on the actual score.

The elementary gap integrals give
$-M_4/6\le J_1\le M_4/2$, $-3J_2-M_5/4\le0$ and $J\le\rho+1$.
The first two follow by integrating, gap by gap,
\[
 \int_a^b(t-a)t\,dt+\tfrac12\int_a^bt^2dt=(b-a)b^2/2,
\]
and its complementary upper inequality; for the negative part of $J_2$
use $\int_u^v(v-x)x^2dx\le\tfrac13\int_u^vx^3dx$.
Consequently a valid constant bound is
\[
 A_2=\tfrac32b(2u+1)+2uM-l/2.
\]
One may take the smaller of $A_2$ and
\[
 A_1=\tfrac32b(2u+1)+\tfrac32\Gamma(R,l)(2u)
                  +\tfrac38u^2(u+1)+\tfrac8{27}u^3-l/2,
\]
where $\Gamma(R,l)$ bounds the nonnegative part of
\[
 -v^2/4+(R+l)v/3-Rl/2+l^3(2R-l)/(12v^2),\qquad l\le v\le R.
\]
This is $v^{-2}\int_l^v(R-x)x(x-l)\,dx$. For each realized negative
magnitude $v$, Tonelli bounds all its gap contributions by this one
integral; multiplying by $v$ and averaging charges $\Gamma\E X_-^3$.
Thus on negative gaps it charges
$3(\rho J_1-J_2)$ by $3\Gamma(\rho-\gamma)/2$; on positive parts use
$t(\rho-t)\le\rho^2/4$, multiply by $(t-a)\ge0$, and integrate
$\sum\tau_{a,b}\int_a^b(t-a)\,dt=J/2$. This gives $3\rho^2J/8$
and includes the positive part of a central gap with $a<0$.
The remaining term uses
$-K_q\le8\rho^3/27+(3b/2-1/2)\rho$, obtained by
$M_5\ge M_4^{3/2}$ and maximizing $\rho y^2/2-y^3/4$.
Thus $A_*:=\min(A_1,A_2)$ is valid when the chosen $\Gamma$ is certified.

Valid bounds on $|K_1|$ are
\begin{align*}
 k_{\rm rad}&=3ub+\frac{R(R+2-3l/2)}{R+1/2}(u-1),\\
 k_{\rm old}&=6u-3u^2/2+3b,\\
 k_R&=Ru/2+\max_{t=l,u}|R/2-3t/2|g+3b.
\end{align*}
For the first, use the radial identity for $K_1+3\rho\beta$ with ratio
$x(x+2-3\rho/2)/(x+1/2)$, convex in $x$ and decreasing in $\rho$.
Its negative part is at most $31/40$, so require its stated endpoint
ceiling to be at least $31/40$. The second uses $R\le6$ and
$|\E[X|X|^3]-3\gamma|\le3\rho$. The third bounds the two fourth
semimoments by $R(\rho\pm\gamma)/2$. Take $k$ to be their minimum.
Also $U\in[-1/2,2M/3-1/2]$, so use
$v_H=(2/5)\max(1/2,2M/3-1/2)$ (or a separately proved smaller bound).
Then $q_*:=c_-l-g/5>0$ is a valid quadratic curvature.

It remains to have a \emph{complete} Fourier bound
\begin{equation}\label{v5entry:fourier-input}
 |E_n|\le K_F+s_F|H|,\quad
 K_F=\tfrac{16}{25}(L_0+K_H+L_2/\sqrt N),\quad s_F=\tfrac{16}{25}L_H.
\end{equation}
The compensated quadratic-hinge inversion has factor $2\sqrt n/\pi$.
The retained triangle transfer uses the parameters $N,M,g,M_6$ and
sixth-order Taylor estimates on $[0,1/2]$. Its conclusion is
\[
 I_0\le(L_0+L_H|H|)/\sqrt n+L_2/n.
\]
The retained complete six-part high-frequency lemma uses the actual
event \eqref{v5entry:event} and gives $\sqrt n I_\infty\le K_H$.
Every damping, floor and window gate in that lemma must hold; one takes
a minimum only of complete six-term bounds. The event proof does not
permit dropping its Gaussian comparator tail. Neither Fourier lemma
uses a score bound or a defect feedback hypothesis.

Completion of the square now proves the all-threshold bound
\begin{equation}\label{v5entry:scalar}
 \mathcal R_c(H)\le
 T:=K_F+c_+A_*+\frac{(c_+k+v_H+s_F)^2}{4q_*},
\end{equation}
provided $A_*\ge0$, $k\ge0$, $q_*>0$ and $c\in[c_-,c_+]$.
This is a scalar inequality, so using a formal coefficient interval
containing $c$ is sound independently of the actual score.
The first numerical obligation is a finite closed cover of
$[1,7361/4000]$ on which the stated Fourier gates and
$T<Nc_-l$ all hold. The endpoint $\rho=1$ is included; at that point
the law is a symmetric two-point law, and all moment ceilings remain
well-defined. Adjacent cell endpoints must both be covered.

\paragraph{Scope of this note.}
The bounded companion checker verifies the support rational reductions,
five exact finite-law moment tests, and a negative control rejecting
radius three with the same envelope. It does not evaluate a whole raw
cover or prove any local comparison predicate. A separate 24-cell
calculation reports $\mathcal R_c/\rho<10$; its source, outward arithmetic,
coverage and complete Fourier premises require their own review before
composition with Lemma~\ref{v5entry:ten}. Subsequent support
contractions, central-gap acquisition, confinement and LOW/HIGH exclusion
still require their own ordered hypotheses. In particular the present
entry theorem alone does not remove the independent finite base.

\section{What the first cover proves}
\label{v5num:section}

The numerical obligation in Section~\ref{v5entry:raw} is small enough to
test directly. Take $N=42$, $R=96/25$ and the 24 closed intervals with
endpoints
\[
 r_i=1+\frac{3361i}{96000},\qquad 0\le i\le24.
\]
Adjacent intervals share their endpoint; no boundary or sign class is
discarded. Enclose the fixed coefficient, using Machin's formula and
rational square-root bounds, by
\[
 c_-=0.4097321837023<\CE<c_+=0.4097321837025.
\]
These numbers do not enclose the maximizing law's actual score.

On $[l,u]$, use the moment bounds of \eqref{v5entry:mom}. The low
triple is the sixth-order triangle estimate; the high estimate is the
minimum of six complete six-term bounds, at windows
$1/8,4/25,9/50,1/5,11/50,1/4$. Every window's gates are checked before
its complete bound is admitted. With the notation of the preceding
section, set
\[
 q(t)=tl-g/5,\quad f(t)=tk+v_H+s_F,\quad
 B(t)=K_F+tA_*+\frac{f(t)^2}{4q(t)}.
\]
For fixed $H$, the pre-completion expression
$K_F+tA_*+f(t)|H|-q(t)H^2$ is affine in $t$.
Its supremum over $H$ is convex. Hence the sharper endpoint choice
\[
 T_i=\max\{B(c_-),B(c_+)\}
\]
is valid for all intermediate $t$, including $\CE$. All divisors are
strictly positive. The retained source functions and the explicit rational
reconstruction are listed in \texttt{cell.md}; \texttt{cell\_check.py}
computes the full cover afresh.

The exact fractions imply the outward rounded bounds
\begin{equation}\label{v5num:raw-results}
 \max_i T_i/r_i<9.455732088645<10,\qquad
 \min_i(42c_-r_i-T_i)>13.996057920754.
\end{equation}
Thus, conditional on the full stated Fourier estimates, the first license
holds on the entire initial moment domain. The last cell has
$T_{23}<17.069861855816$, above the old sufficient uniform cutoff
$2N/5=16.8$. It nevertheless passes the moment-sensitive cutoff.
This example explains why the change of gate matters.

\subsection{Three ordered refinements}
Keep a raw ceiling $T$, a conserved-resource ceiling $U$ and a defect
ceiling $d$ as separate quantities. Initially $U=T$ and
$d=6U/(N\phi_-)$, with $\phi_-=997/2500<\phi(0)$.
All retained ceilings in the following computation are rounded upward
to multiples of $10^{-10}$.

At each step, the previous raw ceiling gives
\[
 \phi_{\min}=\phi_-\left(1-\frac{T}{Nc_-l}\right)>0.
\]
The gap bound $J\le\rho+2m$, where $m=\E X_+\le1/2$, implies
\[
 D\ge\sqrt{10}\rho-\gamma-6m.
\]
It is an inequality, not an identity. Thus, with
$\sigma_-=3.162277660<\sqrt{10}$, the previous defect ceiling yields
\[
 \gamma_{\min}=\max(-g,\sigma_-l-3-d),\qquad
 \rho-\gamma\le\min(2u,(1-\sigma_-)l+3+d).
\]
Use the latter to improve only the negative-gap part of $A_1$ in
Section~\ref{v5entry:raw}; retain its minimum with the old constant bound.
The sign-correct new curvature is
\[
 q(t)=\max\left(tl-g/5,\;
 tl+\tfrac12\min(\phi_{\min}\gamma_{\min},(2/5)\gamma_{\min})\right).
\]
All bounds are frozen during the endpoint maximum. Both alternatives
have slope $l$ in $t$, so the same convexity argument applies.
Recompute the raw upper envelope and take its minimum with the previous
$T$. This new raw bound now licenses the phase price
\[
 \eta=\tfrac12(c_-l-T/N)>0.
\]
Lemma~\ref{v5entry:phase} bounds the conserved resource by the endpoint
envelope with denominator $4(q(t)+\eta)$. Take its minimum with the
previous $U$, then update $d$ by $6U/(N\phi_-)$.
Every update uses already justified bounds; the desired defect is never
an input to its own proof. In particular, $U$ is never substituted for
$T$ in the Gaussian-weight license.

The following are outward upper bounds from exactly three steps, with
all 24 cells retained:
\begin{center}
\begin{tabular}{@{}lrrrr@{}}
\toprule
Stage & maximum $T$ & maximum $U$ & maximum $d$ & cells with $d<4$\\
\midrule
Initial &17.0698618559&17.0698618559&6.1147234045&---\\
Step 1 &17.0698618559&13.9304183554&4.9901197720&19\\
Step 2 &17.0698618559&13.9304183554&4.9901197720&19\\
Step 3 &17.0698618559&13.9304183554&4.9901197720&19\\
\bottomrule
\end{tabular}
\end{center}
The last resource bound also gives the weighted statement
$nD\le209.585030422267$; it is not inferred merely from an unweighted
bound on $D$. Some lower moment cells improve at steps 2 and 3, but the
worst cell does not. In the final cell the acquired lower bound on
$\gamma$ is still $-1.84025$.

The five cells numbered 19--23 remain outside the sufficient $D<4$
regime. A finer bound coupling the moments, or a new structural
restriction, is needed there. The first license has been checked; the
full initial acquisition and the local exclusion have not. This 24-cell
cover must not be compared with the size of the older dense acquisition
as though both computations established the same conclusion.

\section{A smaller cost for count crossings}
\label{v5cc:section}

The whole-count comparison contains a discrete Taylor remainder of the
binomial CDF. Its present bound pays separately for positive and negative
crossings, although those events are disjoint. It also uses a third-moment
estimate for crossings that can occur only beyond a fixed positive distance.
Both losses can be removed without changing the law or the comparison target.

Here is the coupling, including its endpoint convention. Let $U_{0,i}$ and
$U_{1,i}$, $1\le i\le n$, be independent centered residuals, with
$U_{0,i}\le a_0$, $U_{1,i}\ge-a_1$ and
$c=1-a_0-a_1\in(1/2,1]$. Put
\[
 V_j=\sum_{i\le j}U_{1,i}+\sum_{i>j}U_{0,i},\qquad0\le j\le n,
\]
and extend $V_j$ constantly outside $[0,n]$. The sequence $j+V_j$ is
strictly increasing, with increments at least $c$. For a legal count $k$
and $|u|\le1/2$, define
\[
 K=\max\{j\in\mathbb Z:j+V_j\le k+u\},\quad R=K-k,
 \quad t=\E V_k^2.
\]
If an independent selector has law $\operatorname{Bin}(n,p)$, its closed
event occurs exactly when the selector is at most $K$. Thus the mixture CDF
is $\E F_B(k+R)$. This construction uses the actual conditional laws; it
does not replace the convolution by independent choices of moment bounds.

Write $v=p(1-p)$, $\pi_j=\binom njp^j(1-p)^{n-j}$, with zero extension,
$m_B=\max_j\pi_j$, and $\chi=(n+1)p-k$. Assume $\chi<2$. For integer $r$,
the binomial difference identity and double summation give
\begin{equation}\label{v5cc:discrete}
 F_B(k+r)-F_B(k)-r\pi_k\le
 \frac{m_B}{(n+1)v}
 \left\{(\chi-1)_+r_++\chi_+\binom{r_-}{2}
                                      +\binom{r_-}{3}\right\}.
\end{equation}
For completeness,
$\pi_j-\pi_{j-1}=\pi_j((n+1)p-j)/(p(n-j+1))$.
When its numerator is positive, the denominator is at least $(n+1)v$.
For $r>0$, the positive differences have weights $r-i+1$ at $j=k+i$;
only $i=1$ can contribute positively because $\chi<2$.
For $r=-h$, the weights are $h-1-i$ at $j=k-i$, $0\le i\le h-2$.
The bound $(\chi+i)_+\le\chi_++i$ gives the two binomial coefficients
in \eqref{v5cc:discrete}. Zero extension handles counts outside $[0,n]$.

\begin{lemma}[Disjoint costs and a fourth-moment bound]
\label{v5cc:crossings}
Put $d=2c-1/2>0$ and
\[
 M_2(\chi,c,u)=\max\left\{
 \frac{(\chi-1)_+}{(c-u)^2},\,
 \frac{\chi_+}{[c+\min(u,c/3)]^2}\right\}.
\]
Then, pointwise,
\begin{align}
 (\chi-1)_+R_++\chi_+\binom{R_-}{2}
       &\le M_2(\chi,c,u)V_k^2,\\
 \binom{R_-}{3}&\le (V_k)_+^3/d^3,
 \qquad \binom{R_-}{3}\le (V_k)_+^4/d^4.
 \label{v5cc:fourth}
\end{align}
If also $|U_{s,i}|\le A$ and $t\le T$, then
\begin{equation}\label{v5cc:expected}
 \E\{F_B(k+R)-F_B(k)-R\pi_k\}
 \le\frac{m_Bt}{(n+1)v}
 \left[M_2(\chi,c,u)+
 \min\left\{\frac{\sqrt{A^2+3T}}{d^3},
                   \frac{A^2+3T}{d^4}\right\}\right].
\end{equation}
\end{lemma}
\begin{proof}
On $R=r\ge1$, crossing monotonicity gives $-V_k\ge cr-u$.
The difference
$(cr-u)^2-r(c-u)^2=(r-1)(c^2r-u^2)$ is nonnegative.
On $R=-h\le-2$, the count immediately above $K$ is outside the closed
event, so $V_k>c(h-1)+u$. Set $j=h-1\ge1$.
For $u\le c/3$, the function
$j(j+1)/(2(cj+u)^2)$ decreases for $j\ge1$: its derivative has the sign
of $j(2u-c)+u\le0$. Its maximum is $1/(c+u)^2$.
For $u>c/3$, lower $u$ to $c/3$ in this bound.
The positive and negative crossing events are disjoint. Hence the two
quadratic prices use a maximum, not their sum. When $R=0$ or $-1$,
both costs vanish.

For $h\ge3$, put $j=h-1\ge2$. Since $u\ge-1/2$,
\[
 \frac{\binom h3}{[c(h-1)+u]^3}
 \le\frac{j^3-j}{6(cj-1/2)^3}\le\frac1{(2c-1/2)^3}.
\]
For the last inequality the derivative numerator has the sign of
$2cj-(3j^2-1)/2<0$, using $c\le1$ and $j\ge2$; equality of the
last bound occurs at $j=2$. This proves the cubic estimate.
On this same event $V_k>d$, so $(V_k)_+^3/d^3\le(V_k)_+^4/d^4$.
For the other values of $R$, the left side is zero. Thus the quartic
estimate is pointwise too; no approximation at a crossing atom is used.

Independence and centering give, with variances $s_i$ of the $n$ terms,
\[
 \E V_k^4=\sum_i\E Z_i^4+6\sum_{i<j}s_is_j
 \le A^2t+3t^2\le t(A^2+3T).
\]
Cauchy--Schwarz also gives $\E(V_k)_+^3\le t\sqrt{A^2+3T}$.
Take expectations in the three pointwise estimates and then in
\eqref{v5cc:discrete}. The result includes $t=0$.
\end{proof}

The fourth-moment price is smaller whenever
$A^2+3T<(2c-1/2)^2$. It uses a moment bound already present in the
old proof. In that domain the square root is unnecessary for this term.
For larger noise retain the cubic price; taking the minimum preserves
both regimes. Dropping the individual fourth moments $A^2t$ would be
false, even for a centered three-point variable with rare nonzero values.

\subsection{The phase maximum still has three candidates}

In a closed scalar cell, $\chi\le b+u$. Define
\[
 g(u)=\max\left\{
 \frac{(b+u-1)_+}{(c-u)^2},\,
 \frac{(b+u)_+}{[c+\min(u,c/3)]^2}\right\},
 \qquad u_*=c-2b.
\]
For $[l,h]\subseteq[-1/2,1/2]$, its exact maximum is
\begin{equation}\label{v5cc:phase}
 \max\left\{g(l),g(h),
 \left.\frac1{4(c-b)}\;\right|\;
 c>b,\ l\le u_*\le h,\ u_*\le c/3\right\}.
\end{equation}
An ineligible third value is omitted. To prove this, the first summand
inside the maximum is nondecreasing: on its positive region its derivative
has positive numerator $c+u+2b-2>c-u>0$.
Below $c/3$, the positive part of the second has derivative
$(c-2b-u)/(c+u)^3$; above $c/3$ it is nondecreasing.
Its only possible interior maximum is therefore $u_*$, with value
$1/[4(c-b)]$. At that point $b+u_*=(c+u_*)/2\le3/4$, so the first
term vanishes. The supremum of a maximum is the maximum of the two
suprema. This proves \eqref{v5cc:phase}, including singleton intervals,
kinks and $c=1$.

For every allowed parameter, $g$ is bounded by the old sum of these two
terms. Both terms are nonincreasing in $c$. Consequently a conservative
cache at a smaller acquired $c$ is still valid, and the new exact phase
maximum is no larger than either the previous exact sum maximum or its
64-cell outward bound.

\subsection{What this changes in the local comparison}

Use the same acquired amplitude $A$, upper own-count variance $T$ and
lower gap coefficient $c$ as the current cell. Replace its curvature
coefficient by \eqref{v5cc:expected}, maximizing the phase by
\eqref{v5cc:phase}. All other terms, including the marked-crossing linear
term $\pi_k\E R$, the Gaussian repair and the actual binomial mass bound,
are retained. The new coefficient is never larger than the old one.
When passing from $1/n$ to $1/N$, clamp the complete improved inverse-index
correction at zero; its old strict positivity need not survive an improvement.
The existing negative event bracket is still required before using a lower
binomial mass bound.

For a fixed noise band, $A,T,c$ are frozen valid bounds. Hence this
replacement changes only the constant part of the inverse-index correction;
the existing convex endpoint argument in the noise parameter is unchanged.
It does not prove the acquisition of those bounds, cover all count cells,
or replace the separate zero-noise binomial theorem. It supplies a smaller
curvature contribution under exactly the old moment and coupling hypotheses.

\section{One response for the actual moment polynomial}

We give a conditional, uniform improvement on a recorded central probability
panel. It bounds all the moment terms of the fixed-target raw remainder
together. The original gap integrals and the Fourier remainder remain
explicit: the statement does not establish the acquisition of this panel
for a fixed-target maximizer, or a complete local exclusion.

Let $P$ have mean zero, variance one and bounded support, and put
$p=P(X>0)$, $q=1-p$, $m=\E X_+=\E X_-$.
Assume a central support gap $(-A,B)$, with $A,B>0$.
Write $V=\E[X\mid\mathbf1_{X>0}]$, $W=X-V$,
$a=m/q$, $b=m/p$, and $\theta=\E W^2$.
Thus $V=-a$ or $b$, and $m^2=(1-\theta)pq$.
Use the same actual gap quantities $J,J_1,J_2$ as in the conserved identity,
and write $\sigma=\sqrt{10}$, $p_*=(4-\sigma)/2$ and
$D=(3+\sigma)\rho-\gamma-3J$.
For a law $Y$ in the original physical units, denote its moments by
\[
 (\rho_Y,\beta_Y,\gamma_Y,M_Y,E_Y,Z_Y)
 =(\E|Y|^3,\E Y|Y|,\E Y^3,\E Y^4,\E|Y|^5,\E Y|Y|^3).
\]
The unsubscripted moments refer to $X$, and subscript $0$ refers to $V$.
In particular, $V$ has variance $1-\theta$, not necessarily one.

Fix $c=\CE=(3+\sqrt{10})/(6\sqrt{2\pi})$ and
$\varphi=\phi(H/\sqrt n)$, where $\phi$ is the standard normal density
and $H$ is the raw threshold.
For any moment tuple $M$ define the polynomial
\begin{align*}
 T_c(M;J,J_1,J_2,H,\varphi)
 ={}&c\left\{\frac32\beta J+3\rho J_1-3J_2
 -\frac{E}{4}+\frac{\rho M_4}{2}+\frac32\beta\gamma-\frac\rho2\right\}\\
 &+H\left\{c\left(Z-3\beta-\frac32\rho\gamma\right)
                  +\varphi\left(J_1+\frac{M_4-3}{6}\right)\right\}
 -H^2(c\rho+\varphi\gamma/2).
\end{align*}
Thus the actual raw remainder is $\mathcal R_c(H)=T_c(M(P);J(P),J_1(P),
J_2(P),H,\varphi)+E_n(H)$.
Every comparison below keeps the actual $J(P),J_1(P),J_2(P)$ on both sides.
Evaluating this polynomial at $M(V)$ is not evaluating the remainder of
a Bernoulli law.

\subsection{Exact common-measure identity}
On the two sign classes put
\[
 d\mu_+=W^2\mathbf1_{X>0}\,dP,\qquad
 d\mu_-=W^2\mathbf1_{X<0}\,dP.
\]
For magnitude $x>0$ and conditional mean magnitude $v>0$, define
\[
 Q_3=x+2v,\quad Q_4=x^2+2vx+3v^2,\quad
 Q_5=x^3+2vx^2+3v^2x+4v^3.
\]
The identity
$x^j-v^j-jv^{j-1}(x-v)=(x-v)^2Q_j(v,x)$,
together with conditional centering, gives the moment differences.
For example, $\delta\rho=\sum_\pm\int Q_3\,d\mu_\pm$,
$\delta\gamma=\int Q_3\,d\mu_+-\int Q_3\,d\mu_-$,
and $\delta\beta=\mu_+(\R)-\mu_-(\R)$.

For $\epsilon\in\{+1,-1\}$, with $v=b$ or $a$ respectively, put
\begin{align*}
 P_{0,\epsilon}(v,x)
 ={}&-\frac{Q_5}{4}+\frac\rho2Q_4
 +\left(\frac{M_0}{2}-\frac12+3J_1+\frac32\epsilon\beta\right)Q_3
 +\frac32\epsilon(\gamma_0+J),\\
 P_{1,\epsilon}(v,x)
 ={}&\epsilon(Q_4-3)-\frac32(\epsilon\rho+\gamma_0)Q_3,\\
 \mathscr P_\epsilon(v,x)
 ={}&cP_{0,\epsilon}
 +H\left(cP_{1,\epsilon}+\frac\varphi6Q_4\right)
 -H^2(c+\epsilon\varphi/2)Q_3.
\end{align*}
Then the following identity is exact:
\begin{equation}\label{response:exact}
 T_c(M(P);J,J_1,J_2,H,\varphi)-T_c(M(V);J,J_1,J_2,H,\varphi)
 =\sum_{\epsilon=\pm}\int\mathscr P_\epsilon\,d\mu_\epsilon.
\end{equation}
To check the products, use
\[
 \rho M_4-\rho_0M_0=\rho\delta M_4+M_0\delta\rho,
 \quad\beta\gamma-\beta_0\gamma_0=\beta\delta\gamma+\gamma_0\delta\beta,
 \quad\rho\gamma-\rho_0\gamma_0=\rho\delta\gamma+\gamma_0\delta\rho.
\]
The coefficients contain actual moments of the same law. No independent
choice of incompatible moment maxima is asserted to be a law.

\subsection{The cost of the common measures}
Let $\Delta_A=a-A$, $\Delta_B=b-B$ and define
\[
 B_0=\frac{(6+2\sigma)(1-\theta)^{3/2}}{\sqrt{pq}}(p-p_*)^2,
 \quad E_0=6m\{(A+B)(\Delta_A+\Delta_B)+\Delta_A\Delta_B\}.
\]
The weighted gap inequality gives, for $r=D-B_0-E_0$,
\begin{equation}\label{response:budget}
 0\le\sum_\pm\int k_\pm\,d\mu_\pm\le r,
 \quad
 k_+(x)=(\sigma-1)(x+2b)+6B,
 \quad k_-(x)=(\sigma+1)(x+2a)+6A.
\end{equation}
For completeness, the outer-gap estimate implies
\[
 D\ge(\sigma-1)\E X_+^3+(\sigma+1)\E X_-^3
       -6mAB+6BR_++6AR_- ,\quad
 R_+=\E X_+^2-mB,\ R_-=\E X_-^2-mA.
\]
Substitute
$\E X_+^3=mb^2+\int(x+2b)\,d\mu_+$,
$R_+=\mu_+(\R)+m\Delta_B$, and the reflected identities.
The deterministic part is $B_0+E_0$, by
$(\sigma+1)p^2+(\sigma-1)q^2-6pq=(6+2\sigma)(p-p_*)^2$.
This proves \eqref{response:budget}. All costs and both measures are
nonnegative. In particular a pointwise bound
$\mathscr P_\pm\le Kk_\pm$, with $K\ge0$, spends this residual only once.

\section{A uniform bound on the recorded central panel}

The domain is taken from zero-based central row 19 of the retained
confinement table. Its acquisition is a premise; historical success
metadata is not a fresh verification. Assume simultaneously
\begin{gather}
 \frac{209}{500}\le p\le\frac{21}{50},\quad |X|\le\frac{46}{25},\quad
 \theta\le\frac{12657}{10^6},\quad \rho\le\frac{214889}{200000},\label{response:domain}\\
 \frac{72659}{500000}\le\beta\le\frac{174581}{10^6},\qquad
 A\ge\frac{810731}{10^6},\quad B\ge\frac{284079}{250000}.
 \nonumber
\end{gather}
These are closed-domain bounds for one actual law.

\begin{theorem}\label{response:theorem}
Under \eqref{response:domain}, for every integer $n\ge1$ and every real $H$,
\begin{align}
 \mathcal R_c(H)\le{}&T_c(M(V);J(P),J_1(P),J_2(P),H,\varphi)
       +c\left(\frac{221}{200}+\frac9{50}|H|\right)r+E_n(H).
 \label{response:main}
\end{align}
Here $c=\CE$, $\varphi=\phi(H/\sqrt n)$, and $r$ is the exact same-law
residual in \eqref{response:budget}. The conclusion bounds the moment
polynomial; it does not supply a Fourier bound for $E_n$ or replace the
actual gaps by those of $V$.
\end{theorem}

\begin{proof}
We first give rational consequences of the acquired domain:
\begin{gather*}
 a\in[21/25,213/250],\quad b\in[29/25,119/100],\quad m>49/100,\\
 A>81/100,\ B>113/100,\quad 1\le\rho<43/40,\quad
 29/200<\beta<7/40,\\
 3/10<\gamma_0<1/3,\quad M_0<10/9,\quad
 J<83/40,\quad \tfrac32(\gamma_0+J)>3,\quad J_1<1.
\end{gather*}
The conditional mean bounds follow by squaring
$a^2=(1-\theta)p/q$ and $b^2=(1-\theta)q/p$.
For $v=pq$, its interval is
$[60819/250000,609/2500]$.
The identities
\[
 \gamma_0=(1-\theta)^{3/2}(1-2p)/\sqrt v,
 \qquad M_0=(1-\theta)^2(1/v-3)
\]
give the two skeleton bounds; squaring positive quantities suffices.
Moreover $J\le\rho+2m\le\rho+1$, whereas $J\ge m(A+B)^2$.
For the last bound, negative-threshold portions of $J_1$ are nonpositive.
On positive outer gaps, $t-u\le t$; on the central positive portion
$t-u=t+A$. Consequently
\begin{align*}
 J_1&\le\int_0^\infty t^2\tau(t)\,dt+\frac{AmB^2}{2}
 =\frac{\E X_+^4}{3}+\frac{AmB^2}{2}\\
 &<\frac{(46/25)(43/40)}3+\frac{(213/250)(119/100)^2}{4}
 =\frac{28828879}{30000000}<1.
\end{align*}
The positive integrands justify summing countably many gaps. This uses the
acquired support before the fourth moment is used.

Since $\sigma>79/25$, define lower costs
\[
 d_+(x)=\frac{14739}{1250}+\frac{54}{25}x,\qquad
 d_-(x)=\frac{14811}{1250}+\frac{104}{25}x.
\]
They are strictly positive and at most $k_+(x),k_-(x)$ respectively.
The constant responses are bounded by
\begin{align*}
 P_{0,+}&\le p_+(x)=\frac{435879881}{36000000}
              +\frac{320027}{90000}x-\frac{17}{400}x^2-\frac14x^3,\\
 P_{0,-}&\le p_-(x)=\frac{895536247}{375000000}
              +\frac{289067}{90000}x+\frac{47}{400}x^2-\frac14x^3.
\end{align*}
Indeed use $\rho/2\le43/80$ and
$M_0/2-1/2+3J_1\le55/18$.
On the positive class use $3\beta/2\le21/80$ and
$3(\gamma_0+J)/2\le289/80$; on the negative class use
$-3\beta/2\le-87/400$ and $-3(\gamma_0+J)/2\le-3$.
For the resulting cubic, the $x^2$ and $x$ coefficients decrease with
the mean magnitude, and the constant coefficient increases. Inserting
the corresponding endpoints gives exactly these two polynomials.

There is a finite, exact check of their ratio maxima. For a cubic $p$
and affine $d=d_0+d_1x$, its derivative numerator is
\[
 G(x)=p_1d_0-p_0d_1+2p_2d_0x+(p_2d_1+3p_3d_0)x^2+2p_3d_1x^3.
\]
For each class, $G''<0$ and $G'$ is negative at the left endpoint;
thus $G$ strictly decreases. Its signs at the two endpoints are opposite,
so precisely one stationary maximum must be checked. Rational enclosing
brackets and strict bounds are
\[
\begin{array}{c|c|c|c}
 &x\text{ domain}&x\text{ at the stationary maximum}&p(x)/d(x)\text{ upper bound}\\\hline
 +&[113/100,46/25]&[1194059,1194060]/10^6&221/200\\
 -&[81/100,46/25]&[1662099,1662100]/10^6&3/8.
\end{array}
\]
The endpoint values are smaller as well. These are rational inequalities:
on a positive interval $[l,u]$, bound each monomial $p_jx^j$ using $u$
when $p_j\ge0$ and $l$ otherwise, then divide by $d(l)>0$.
The accompanying checker prints all exact coefficients, derivative signs,
brackets and endpoint comparisons; no sampled positivity is used.

For the linear term in $H$, use
$0\le\varphi/c\le6/(3+\sqrt{10})<1$.
After division by $c|H|$, the following four polynomials majorize its
response, on the same respective magnitude domains:
\[
\begin{array}{c|c|c}
\operatorname{sgn}H&\operatorname{sgn}X&\text{upper polynomial }s(x)\\\hline
 +&+&-53693/20000+(62/75)x+(7/6)x^2\\
 +&-&3189/1000-(19/80)x-(5/6)x^2\\
 -&+&19321/5000-(83/400)x-x^2\\
 -&-&-157893/62500+(88/125)x+x^2.
\end{array}
\]
For example, the first is obtained from
$(7/6)Q_4-3-(39/20)Q_3$, maximizing over $b\le119/100$.
The others follow from
$-(5/6)Q_4+3+(93/80)Q_3$,
$-Q_4+3+(169/80)Q_3$, and $Q_4-3-Q_3$,
respectively. The mean-endpoint directions follow by differentiating
these displayed quadratics with respect to their mean magnitude.
For each row, $(9/50)d-s>0$ throughout the closed interval: if its
quadratic coefficient is negative, check both endpoints; otherwise
its derivative is positive from the left endpoint, where its value
is positive. All these endpoint tests are exact rational inequalities.

Finally $\varphi<c$ implies $c>\varphi/2$, so the $H^2$ term of each
$\mathscr P_\epsilon$ is nonpositive. We have proved pointwise
\[
 \mathscr P_\pm(v,x)\le c\left(221/200+(9/50)|H|\right)k_\pm(x).
\]
Integrate once against the two positive measures and use
\eqref{response:exact}--\eqref{response:budget}.
\end{proof}

\section{What improves, and what remains}
The gain is more than a new notation for the positive-measure lemma.
At $H=0$, consider the explicit split-response bound that discards
$-Q_5/4\le0$ before taking the positive-class supremum, and uses the same
upper bounds for all remaining coefficients. Its required constant is
at least its value at $x=46/25$, $b=119/100$:
\[
 K_{\rm split}\ge
 \frac{(43/80)Q_4+(55/18+21/80)Q_3+289/80}{d_+(46/25)}
 =\frac{173316497}{113512320}.
\]
In contrast, the combined constant is $221/200$, smaller by at least
$239426917/567561600>0.4218$. For $r>0$ this gives a strict improvement
of that uniform upper bound. This benchmark is an explicitly defined
relaxation, not a claim that any law realizes its independent moment
endpoints, or that this beats every previously implemented atlas bound.

The domain has actual laws with positive noise. For example the masses
\[
 (29/100,29/100,21/100,21/100)\quad\text{at}\quad
 (-533,-517,717,733)/617
\]
give mean zero, variance one, and conditional means $-525/617,725/617$.
Their moments and support satisfy \eqref{response:domain}, with
$0<\theta<12657/10^6$; exact checks are recorded. This law is only a
non-vacuity and identity control, not a claimed maximizing witness.

To turn \eqref{response:main} into a complete raw bound, one must still
control the baseline polynomial with the original $J(P),J_1(P),J_2(P)$
and the full signed Fourier term $E_n(H)$. Conditional mean replacement
changes the gaps, and $V$ is not standardized when $\theta>0$.
Consequently neither the Bernoulli theorem nor a putative identity
$J(P)=J(V)$ supplies these missing steps. The independent conditional
expectation response also does not remove the quadratic remainder in
$f^n-g^n-ng^{n-1}(f-g)$ or provide high-frequency decay.
Finally, the old table concerns its original attained-ratio acquisition;
using its domain for a fixed-target maximizer requires the separately
stated same-law acquisition/transfer hypotheses. No certificate or cell
is removed by this result.


\paragraph{Research and verification record.}
The accompanying source files contain exact checkers, source locations,
input hashes, failed controls and fresh execution records. Independent
review rederived the count-cost conclusion without importing its
constructor. The coordinating review checked the entry and response
derivations and the order of the numerical implications. These are scoped
mathematical reviews, not a formalization or an end-to-end certificate
acceptance. Existing certificates and their original targets are preserved.

\paragraph{Use of AI.}
AI systems were used extensively in developing and checking the arguments,
writing the programs, and drafting this note. The proofs and limitations
are presented for author and independent mathematical review; agreement
between AI systems is not evidence beyond the arguments and checks shown.
\end{document}
